% SOURCE docs/release/2026-10-29/source/docs/writing-repair/2026-10-29/src/ad84e7f88d7e93fe.tex CHAR 14999
\chapter{A learned graph can leave the pricing decision unresolved}
\label{rewrite:discovery}

Suppose $X$ measures a proposed rate change in percentage points, $M$ an
intermediate market quantity and $Y$ gross funding contribution in thousands
of pounds. The units and coefficients in this section are authored for the
example. Consider the centered Gaussian observational covariance
\[
 \Sigma=\begin{pmatrix}
 1&0.6&0.48\\0.6&1.36&1.088\\0.48&1.088&1.8704
 \end{pmatrix}.
\]
The law satisfies $X\perp Y\mid M$. A chain $X\to M\to Y$, a
fork $X\leftarrow M\to Y$, and a reverse chain $X\leftarrow M\leftarrow Y$
can all reproduce it using independent Gaussian errors.

For the first chain, take $M=0.6X+\varepsilon_M$ and
$Y=0.8M+\varepsilon_Y$, with variances $\Var(X)=1$,
$\Var(\varepsilon_M)=1$ and $\Var(\varepsilon_Y)=1$. The intervention
slope of $Y$ on $X$ is $0.6(0.8)=0.48$. For the fork, take
$\Var(M)=1.36$, $X=(0.6/1.36)M+\eta_X$, and $Y=0.8M+\eta_Y$,
with residual variances $1-0.6^2/1.36$ and $1$. The resulting covariance
is the same, but intervening on $X$ has zero effect on $Y$. The reverse
chain uses the corresponding conditional regression coefficients and also
has zero effect. The numerical check reconstructs all three covariance
matrices and verifies positive residual variances.

\section{Keep the causal ambiguity in the financial report}
Within this specified three-model family, the slope set is $\{0,0.48\}$.
At a quarter-point intervention and a stipulated \pounds60 action cost,
net value is
\[
 1000(0.25)\{0,0.48\}-60=\{-60,60\}
 \quad\hbox{pounds}.
\]
There is no sign-robust recommendation across the retained explanations.
Returning the slope from the forward chain alone would hide the unresolved
orientation that determines the action. A probability for each graph would
require additional justified modelling; equal probabilities are not supplied
by observational equivalence. Nor should the discrete set be called a
confidence interval. It represents identification ambiguity even if the
covariance is known exactly.

The example's finite graph family and linear Gaussian law are restrictions.
They do not describe every model compatible with arbitrary banking records.
When the covariance and graph class are estimated, uncertainty must cover
both the observed law and the retained models before mapping them to effects.
An edge bootstrap conditional on one chosen graph does not supply that
combined guarantee. The foundations and detailed temporal-discovery
comparisons retain the arguments for unresolved orientations, latent causes
and class-to-effect propagation.

For a bank, useful additional information could be a defensible ordering
that excludes two orientations, a valid instrument, or an intervention on
$X$. The cost of obtaining that information should be compared with the
consequence of a mistaken action. More precise estimation of this same
covariance cannot distinguish these three laws. The proposed service should
therefore preserve a class or explicit refusal and ask what would narrow it,
instead of completing arrows solely to make a diagram look decisive.

\section{From a known covariance to estimated effect sets}
Chapter~\ref{exec13:temporal-region} supplies the next inferential step under
a restricted stationary Gaussian time-series model. It constructs a region
jointly over the observed covariance, a continuous unknown common lag and a
declared family of causal models, then maps that region to effects. The
coverage argument keeps the true model and observed-law uncertainty together;
it does not condition on an orientation chosen for convenience. The chapter
also shows what a direct intervention can add.

This is a developed restricted construction, not an executed general
bank-discovery pipeline. The maintained graph family, scalar lag and
innovation law must be defensible before applying it. An unrestricted latent
model may leave the whole line admissible. In the opening pricing example,
the narrower set already contains opposite net values, so a more elaborate
estimator is useful only if its information or restrictions can resolve a
consequential ambiguity. A labelled whole-line result can be the correct
answer to a well-defined question. A separate computational issue arises
when a search has not tested the conditioning sets needed to remove a link.
Section~\ref{repair29:discovery} compares a repair on exact-oracle cases and
shows why an unfinished search needs an explicit status. Search completeness,
statistical uncertainty and identification ambiguity are different limitations.


